\chapter{Anexos}

\section{Archivos principales}

\begin{table}[H]
\small
\setlength{\tabcolsep}{3pt}
\begin{tabularx}{\textwidth}{@{}P{0.43\textwidth}Y@{}}
\toprule
\textbf{Archivo} & \textbf{Función}\\
\midrule
\path{coppeliasim/Sim_T2_Phase1_Basic.ttt} & Escena final validada.\\
\path{coppeliasim/Sim_T2_Phase1_Basic_validation.json} & Resultado de validación automática.\\
\path{coppeliasim/build_phase1_basic_scene.py} & Generador de escena y validador.\\
\path{coppeliasim/phase1_r1_task_controller.lua} & Controlador principal de R1: tareas, navegación, SLAM, batería y telemetría.\\
\path{coppeliasim/phase1_b1_walking_manager.lua} & Movimiento y acciones visuales de B1 (Bill).\\
\path{coppeliasim/capture_coppeliasim_screenshots.lua} & Capturas 3D renderizadas desde CoppeliaSim para el informe.\\
\path{coppeliasim/Sim_T2_Phase2_SLAM_Unknown.ttt} & Escena con mapa desconocido, obstáculos no modelados y tarima móvil.\\
\path{coppeliasim/build_phase2_unknown_slam_scene.py} & Generador y validador de Phase 2.\\
\path{coppeliasim/Actividad2_1_Basic_Follower.ttt} & Escena básica de seguimiento de Bill.\\
\path{coppeliasim/pioneer_pid_follower_controller.lua} & Controlador del seguidor con P, PI, PID, LQR y NMPC por disparo directo.\\
\bottomrule
\end{tabularx}
\caption{Escenas y controladores principales.}
\end{table}

\begin{table}[H]
\small
\setlength{\tabcolsep}{3pt}
\begin{tabularx}{\textwidth}{@{}P{0.43\textwidth}Y@{}}
\toprule
\textbf{Archivo} & \textbf{Función}\\
\midrule
\path{coppeliasim/run_phase1_slam_export.py} & Exportación CSV de la ejecución base.\\
\path{coppeliasim/run_phase1_slam_comparison.py} & Comparación de algoritmos SLAM implementados en Lua.\\
\path{coppeliasim/plot_phase1_slam_comparison.py} & Generación de gráficas de comparación SLAM.\\
\path{coppeliasim/phase1_slam_logs/} & CSV, JSON, Markdown y figuras de la comparación SLAM.\\
\path{coppeliasim/run_phase2_slam_unknown_comparison.py} & Exportación comparativa Phase 2 para los cuatro modos SLAM.\\
\path{coppeliasim/plot_phase2_slam_unknown_comparison.py} & Figuras y métricas de reconstrucción Phase 2.\\
\path{coppeliasim/phase2_slam_logs/} & CSV, JSON, Markdown y figuras de Phase 2.\\
\path{coppeliasim/phase1_controller_logs/} & Registros CSV y gráficas de comparación de controladores.\\
\path{coppeliasim/follower_pid_logs/} & CSV y gráficas de seguimiento de Bill.\\
\bottomrule
\end{tabularx}
\caption{Exportadores, registros y escena complementaria.}
\end{table}

\clearpage
\section{Parámetros relevantes}

\begin{table}[H]
\small
\setlength{\tabcolsep}{3pt}
\begin{tabularx}{\textwidth}{@{}P{0.30\textwidth}P{0.18\textwidth}Y@{}}
\toprule
\textbf{Parámetro} & \textbf{Valor} & \textbf{Comentario}\\
\midrule
Radio de rueda & \SI{0.0975}{m} & Conversión de velocidad lineal a angular de motor.\\
Distancia entre ruedas & \SI{0.331}{m} & Modelo diferencial del Pioneer P3DX.\\
Velocidad máxima & \SI{0.55}{m/s} & Saturación de avance.\\
Velocidad angular máxima & \SI{1.32}{rad/s} & Saturación de giro.\\
Rango de evitación & \SI{0.62}{m} & Umbral de influencia reactiva de obstáculos.\\
Parada crítica & \SI{0.12}{m} & Distancia frontal mínima antes de detener el avance.\\
Batería inicial & 96\% & Valor usado en validaciones.\\
Umbral batería baja & 56\% & Limita velocidad si se alcanza.\\
Umbral batería crítica & 26\% & Modo conservador.\\
Resolución de grid & \SI{0.18}{m} & Usada por GMapping, Hector y Cartographer en grid.\\
Margen del planificador & \SI{0.40}{m} & Separación mínima deseada respecto al mapa.\\
Desplazamiento lateral & \SI{0.52}{m} & Distancia lateral del punto intermedio local.\\
Retardo de recuperación & \SI{10.0}{s} & Activa \texttt{RECOVERY\_DIRECT} si no hay progreso hacia el objetivo.\\
Duración de recuperación & \SI{5.5}{s} & Ventana en la que se suspende el mapa para planificar, sin apagar la evitación local.\\
\bottomrule
\end{tabularx}
\end{table}

\section{Fragmentos Lua propios}

\begin{lstlisting}[caption={Conversión de velocidad diferencial a ruedas.},label={lst:wheel-conversion-appendix}]
function DifferentialDrive:setVelocity(v, w)
    local left = (v - 0.5 * self.cfg.trackWidth * w) / self.cfg.wheelRadius
    local right = (v + 0.5 * self.cfg.trackWidth * w) / self.cfg.wheelRadius
    sim.setJointTargetVelocity(self.handles.leftMotor, left)
    sim.setJointTargetVelocity(self.handles.rightMotor, right)
end
\end{lstlisting}

\begin{lstlisting}[caption={Entrega, espera de trabajo humano y retorno de herramienta.},label={lst:handoff-appendix}]
local function deliverToolToBill(tool, workSurface, completedCount, nextSpec, nextState, afterDeliver)
    stopRobot('MANIPULATING')
    if delayElapsed(cfg.manipulationDelay) then
        carryingTool = 0
        placeToolAt(tool, workSurface, 0.43)
        completedTaskCount = math.max(completedTaskCount, completedCount)
        currentTaskSpec = nextSpec
        setShapeColorSafe(tool, ToolColor.onTable)
        if afterDeliver then afterDeliver() end
        setTaskState(nextState)
    end
end

local function waitForBillWork(tool, workSurface, dropPoint, workDelay, nextTaskId, nextSpec, nextState)
    stopRobot('WAIT_B1')
    placeToolAt(tool, workSurface, 0.43)
    if delayElapsed(workDelay) then
        currentTaskId = nextTaskId
        currentTaskSpec = nextSpec
        placeToolAt(tool, dropPoint, 0.36)
        setTaskState(nextState)
    end
end
\end{lstlisting}

\begin{lstlisting}[caption={Animación de Bill trabajando con brazos hacia la mesa.},label={lst:bill-arms-forward}]
elseif self:actionContains(action, 'WORKING') then
    pose.leftShoulder = 0.62 + 0.06 * pulse
    pose.rightShoulder = 0.60 - 0.06 * pulse
    pose.leftElbow = 0.28 + 0.05 * math.abs(tap)
    pose.rightElbow = 0.26 + 0.05 * math.abs(tap)
    pose.leftKnee = 0.07 + 0.03 * math.abs(pulse)
    pose.rightKnee = pose.leftKnee
\end{lstlisting}

\section{Comparación de control}

\begin{table}[H]
\footnotesize
\setlength{\tabcolsep}{3pt}
\begin{tabular}{@{}P{0.12\textwidth}P{0.13\textwidth}P{0.13\textwidth}P{0.13\textwidth}P{0.13\textwidth}P{0.14\textwidth}P{0.11\textwidth}@{}}
\toprule
\textbf{Modo} & \textbf{Dur. [s]} & \textbf{Mov. [s]} & \textbf{Ruta [m]} & \textbf{Bat. [\%]} & \textbf{Riesgo max.} & \textbf{Puntaje}\\
\midrule
P & 115,20 & 37,70 & 12,843 & 19,804 & 0,410 & 0,573\\
PI & \textbf{110,30} & \textbf{31,20} & 12,867 & 19,516 & 0,487 & \textbf{0,274}\\
PID & 110,85 & 31,55 & 12,856 & 19,784 & 0,469 & 0,567\\
LQR & 114,15 & 36,20 & 12,796 & 19,505 & \textbf{0,401} & 0,436\\
NMPC & 112,30 & 33,05 & \textbf{12,788} & \textbf{19,316} & 0,444 & 0,326\\
\bottomrule
\end{tabular}
\end{table}

\begin{figure}[H]
  \centering
  \includegraphics[width=\textwidth]{phase1_controller_metrics.png}
  \caption{Métricas complementarias de control: tiempos, trayectoria, suavidad, batería, riesgo y error de pose.}
  \label{fig:controller-metrics}
\end{figure}

\section{Notas de alcance}

\begin{itemize}
  \item Las variantes GMapping, Hector y Cartographer son implementaciones Lua
  inspiradas en esos métodos y calculadas con los datos de los sensores de la escena.
  \item El sensor se representa con los 16 ultrasonidos del Pioneer visualizados
  como rayos. Las figuras de reconstrucción SLAM corresponden a esos rayos de
  proximidad y al arreglo ultrasónico simulado.
  \item \texttt{phase2MappingEvidencePct} mide actividad de mapeo. La IoU y la
  cobertura geométrica contra una referencia real quedan fuera de esa señal.
  \item La planificación local por punto intermedio resuelve bloqueos de una celda
  pequeña sin introducir una capa global adicional.
  \item LQR y NMPC se usan como comparaciones acotadas; el análisis central usa
  métodos clásicos de seguimiento, evitación y estimación.
\end{itemize}
